% s-234-p234-m001; original PDF 234, 235
\Aphorism{s-234-p234-m001}{328}
\Sperr{„}Was meinen Zarathustra anbetrifft, schreibt Nietzsche, so lasse ich Niemanden als dessen Kenner gelten, den nicht jedes seiner Worte irgendwann einmal tief verwundet und irgendwann einmal tief entzückt hat“. Ja, das ist es! tief verwunden und tief entzücken, nur nicht in der gleichgültigen Mitte der Zu- und Abneigung wandeln, nur nicht ungeliebt und unbefehdet „klassisch“ sein! Sehen wir uns die wirklich seelenbewegenden Menschen und Kunstwerke an: alle entfachen den flammendsten Streit um Für und Wider, nicht nur in des einzelnen Menschen Empfinden, wenn sie darin auf ein verschiedenes „Irgendwann einmal“ treffen, sondern mehr noch im Empfinden vieler Mit- und Nachlebenden, selbst wenn man die feinsten und tüchtigsten heraussondert und die abhängigen bei Seite lässt. Jedesmal geht ein neues Staunen durch die Welt über so viel Kluft und Zwiespalt im künstlerischen Urtheil, jedesmal verwundert sich das eine Lager über all die Blindheit und all den bösen Willen im feindlichen; und doch ist nichts natürlicher, als dass ein Kunstwerk, das überhaupt die Herzen aus der Brust heraufholt, bei Jedem Verschiedenes zu Tage bringt. Es ist nicht, wie die Parteien selbst meinen, die Sonderung in Vorgeschrittene und Rückständige, in Erlauchte und Niedrige, in Adel und Pöbel, die sich hier vollzieht — davon sehe ich ab, und das Pressgesindel, das etwa einen Richard Wagner umkläfft hat, wird Niemand zu seinen Gegnern rechnen; es ist einfach die Sonderung in Individuen, deren jedes, aus anderer Heimath stammend, an anderer Bildung gereift, nach anderen undiscutirbaren Gründen wählt und richtet. Die grenzenlos bezaubernde Wirkung auf mich wird durch die grenzenlos abstossende auf dich erkauft, unbewundert und unbestritten bleibt allein das Mittelmässige. Goethe’s Wahlverwandtschaften könnten nicht dem Einen die reifste und intimste Liebesdichtung sein, wenn sie nicht dem Anderen ein greisenhaftes Problem in geschraubter Behandlung wären; Tristans und Isoldens schwelgende Ekstase muss dem Einen Pein und Scham verursachen, damit sie den Anderen in alle sieben Himmel entzücken kann. Über Haydn giebt es keine Meinungsverschiedenheit; der ist „klassisch“!
\par

